\section{Discussion}\label{sec:discussion}

\textbf{Repeated runs support a descriptive ranking, with substantial evaluation qualifications.} Dynamic-DeepHit has the highest recorded mean C-index and mean time-dependent AUC in all three scenarios. Hazard Transformer ranks second and performs substantially better on the full twenty-feature cohort than in the feasibility pilot. Most other learned models have discrimination near KFRE in the four-/eight-feature cohorts. These results describe the implemented models under this evaluation pipeline; they do not establish statistical equivalence, clinical superiority, or the intrinsic superiority of an architecture.

\textbf{Feature count is confounded with cohort and representation.} The eight-feature complete-case cohort is smaller than the four-feature cohort, and the twenty-feature representation includes missingness indicators and a much broader patient group. Dynamic-DeepHit and Hazard Transformer use patient histories, whereas several comparators use final observations. Both population selection and access to longitudinal information can affect performance. A matched-patient, matched-history ablation is needed to isolate the contribution of additional measurements or architecture.

\textbf{Discrimination and absolute risk require separate assessment.} The transformed-score IBS favors Dynamic-DeepHit, but its native predictions can still substantially underpredict risk, as in twenty-feature rep1. Hazard Transformer's constant two-year probabilities further show why a favorable scalar-score ranking does not establish clinically useful risk estimates. Because the current decision-curve comparators use different evaluation units, the completed reports cannot support recommendations about model-guided referral.

\subsection{Limitations}\label{sec:limitations}
\begin{itemize}
\item \textbf{One data source and a fixed patient partition.} Five model runs quantify fitting/tuning variation, not uncertainty from independent patient samples. There are no patient-bootstrap confidence intervals, paired significance tests, or external-cohort validation. Deterministic results can have zero displayed SD without being known precisely. The selected hospital cohort is event-rich and differs from the populations supporting KFRE's published validation; performance here should not be generalized to routine clinical screening.
\item \textbf{Evaluation and model selection require further validation.} Although test predictions are patient-level, the reported IPCW metrics use raw training rows to estimate the censoring distribution. IBS evaluates the common transformed risk scores rather than each model's native survival distribution. Logistic Hazard uses the nominal test set for tuning and checkpoint selection, so its numbers are not an untouched test-set estimate. These pipeline choices limit comparative interpretation and require a harmonized validation rerun.
\item \textbf{Prediction time and available information are not fully aligned.} Final-observation scoring for some models and full-history scoring for others are not a common prospective landmark design. Whole-history nearest-uACR matching and bidirectional chemistry matching can attach measurements taken after an anchor. Quantified on the extracted cohort, the whole-history uACR match attaches a value recorded \emph{after} the anchor creatinine for 40.4\% of four-feature rows (40.1\% of eight-feature rows), with a median absolute anchor-to-match separation of 264 days and a 95th percentile of 5.4 years; the $\pm$24h chemistry-panel matches are by contrast simultaneous with the anchor for 92.7--98.6\% of rows (serum albumin, the sparsest, is simultaneous for 54.6\% and later for 18.9\%, and is matched at all for only 39.8\% of anchor rows). The follow-up available before the event is correspondingly short: for label-positive patients the median time from their first extracted laboratory record to the labelled event is 1.2 days in the four-feature cohort, 0.4 days in the eight-feature cohort and 104 days in the twenty-feature cohort, and 18--35\% of patients contribute a single exported row, for whom the time-varying input reduces to one snapshot. A two-year horizon measured from that origin is a follow-up window, not a prediction lead time. The current results therefore do not establish performance using only information available at a prespecified clinical decision time.
\item \textbf{The labelled outcome is not the KFRE endpoint.} The event is an acute or unspecified kidney-failure diagnosis code recorded in a patient with CKD stage 3--5, dated to the admission that carried it, not the initiation of maintenance dialysis or preemptive transplantation that KFRE predicts (Section~\ref{sec:outcome}). Reported calibration differences between KFRE and the fitted models therefore combine an endpoint difference, a population difference, and the fitted/externally-specified asymmetry, and this benchmark does not separate them.
\item \textbf{Below-chance and degenerate predictions remain unresolved.} Cox has mean C-index 0.308, 0.470, and 0.457 across the three scenarios, and eight-feature DeepSurv averages 0.484. Hazard Transformer has constant or near-constant native two-year probabilities. Repetition of these behaviors does not exclude a shared implementation or representation issue, and no causal explanation is established here.
\item \textbf{No competing-risk assessment or subgroup validation.} Death before ESRD is not explicitly modeled as a competing event. The benchmark also does not establish performance across demographic subgroups or care settings.
\item \textbf{Importance diagnostics are heterogeneous and partly incomplete.} Coefficients, gradients, and permutation scores have different meanings and scales. All twenty-feature Dynamic-DeepHit and Hazard Transformer gradient analyses failed memory allocation; their fallback outputs do not support substantive interpretation.
\end{itemize}

\subsection{Future Work}
The next evaluation should reserve a separate tuning partition, define a common prospective landmark and information cutoff, estimate censoring from patient-level training outcomes, and compare native survival probabilities and calibration under a common protocol. Decision curves require the same patients and thresholds for every strategy. Matched-cohort feature ablations, patient-level uncertainty estimates, external validation, explicit competing-risk analysis, and memory-bounded feature-importance calculations would then address the remaining gaps.
